


One of the primary works used for conceptualizing the theory of ontological security in regard to states is Jennifer Mitzen's \parencite*{mitzenOntologicalSecurityWorld2006} 2006 article titled \textit{Ontological Security in World Politics: State Identity and the Security Dilemma}. This article expands upon Giddens' \parencite*{giddensModernitySelfidentitySelf2009} work focusing on how individual identities are constructed, and she applies the concept to state actors as a direct response realist theorists \Parencite[][p. 342-344]{mitzenOntologicalSecurityWorld2006}. Mitzen's \parencite*{mitzenOntologicalSecurityWorld2006} application of ontological security is largely covered in the following theoretical framework section, but it created a framework that has been cited in much of the literature on this topic, including much of the literature cited in this project.

There have been a few works that focus on ontological security on BiH, primarily Kočan's book \textit{Identity, Ontological Security and Europeanisation in Republika Srpska} \parencite*{kocanIdentityOntologicalSecurity2023}, Lea David's article ``Policing Memory in Bosnia'' \parencite*{davidPolicingMemoryBosnia2019}, and Jagiełło-Szostak's ``Srebrenica as a Marker of Memory in Bilateral Relations Between Serbia and Bosnia and Herzegovina in the Light of Ontological Security'' \parencite*{jagiello-szostakSrebrenicaMarkerMemory2025}. A foundation of all of these works is a focus on the memory of BiH and how the Yugoslav Wars, in particular, shape the ontological security of BiH, regardless of ethnic divides \parencites[][p. 107-119]{kocanIdentityOntologicalSecurity2023}[][p. 216-220]{davidPolicingMemoryBosnia2019}[][p. 4-5]{jagiello-szostakSrebrenicaMarkerMemory2025}. Despite many similarities in the theoretical approaches to Bosnian identity and ontological security, the subjects of them differ. Kočan \parencite*[][p. 120-125]{kocanIdentityOntologicalSecurity2023} focuses on Bosnian Serbs, particularly in Srbska Republika, which provides insight into this particular enclave, which constitutes the second largest enclave, that also has a strong impact on foreign relations due to their ties with Serbia. As well focusing directly on Srbska Republika, he also elaborates on how the EU's attempts of \textit{Europeanisation} have strained the ontological security in BiH and internal relations between parties \Parencite[][p. 125-130]{kocanIdentityOntologicalSecurity2023} As well as substantial insight into the context regarding identity and ontological security of Republika Srpska, Kočan's \parencite*[][p. 149-155]{kocanIdentityOntologicalSecurity2023} methodological approach follows a similar intended methodology as this project, revolving around content and discourse analysis, providing a strong example of utilizing mixed methods of text analysis for analysis of ontological security and identity. David \parencite*[][p. 216-220]{davidPolicingMemoryBosnia2019} and Jagiełło-Szostak \parencite*[][p. 6-10]{jagiello-szostakSrebrenicaMarkerMemory2025} put more emphasis on memory in BiH, specifically how it is controlled and utilised, where Daivd focuses on mnemonic politics as a way of shaping indentities and ontologies, while Jagiełło-Szostak focuses specifically on the memory of the Srebrenica genocide and the roll it plays on the collective memory, especially between the different ethnicities of BiH. Both approaches provide novel ways of working with the fragile and often contentious memories that stem from Balkan history, and particularly the events of the Second World War and the Yugoslav wars.

Due to the complexity of identity, particularly in the Balkans, further context and history was provided by a selection of works, the first of which is Todorova's \textit{Imagining the Balkans} \parencite*{todorova2009}. Todorova \parencite*[][p. 18-19; 38-41]{todorova2009} largely expands upon concepts such as \textit{orientalism}, as coined by Said \parencite*{said1979}, but does so in a Balkan context that emphasises the history of the region and the perception of the Balkans internally and externally. Her analysis targets the divisions within the understanding of the individual identities throughout the Balkans while also highlighting the \textit{othering} carried out by ``Western Europe'' on the region as a whole, which often exoticizes the region and its peoples through terms such as \textit{balkanization} or \textit{balkanized} \Parencite[][p. 33-35; 39-41]{todorova2009}. Another text that is highly influential in understanding the history of the Balkans, and particulalrly BiH is Andrić's novel \textit{Na Drini Ćuprija (The Bridge on the Drina)} \parencite*{andric1977}. \textit{The Bridge on the Drina} discusses the life in Višegrad in the 16th century, but it provides important insight into the mix of identities in BiH and its history \parencite*{andric1977}. As well as historical context of BiH and the region, the ethnic makeup of BiH also provides further detail that is of importance to this project. Bringa \parencite*{Bringa_2000} provides insight from an anthropoligic perspective regarding the mixed identities in BiH and the complex nature Bosnian identity. Bringa's article provides first hand experiences and observations from mixed villages in BiH, and she describes many of the practices that allow for cohabitation, but still distinction between the different ethnicities throughout BiH \Parencite[][p. 82-84]{Bringa_2000}. BiH has a complex history, which leads to a similarly complex mix of identities, which also provides a basis for this project, which focuses on politicians' attempts of both leveraging the differences of these identities for particular reasons and attempts to homogenize these identities into a common Bosnian identity as well.

As well as historical and contextual literature, this project uses many recent developments in methodological approaches, primarily centered around machine automated text analysis. Natural language processing has largely increase in its capabilities for text analysis in the past 5-10 years, and is increasingly taking a foothold in political science \Parencites{laurer2024}{lucas2015}. Laurer \parencite*[][p. 85-86]{laurer2024} has advocated for the use of supervised transformers based models as a way of increase the amount of data that can be processed for academic text analysis, particularly in social sciences. Methodological apporaches such as deep transfer learning through models such as BERT allow for specialization of parameters based on a particular corpus, and it has the possibility to save time and costs over traditional text analysis methods \Parencite[][p. 87-90]{laurer2024}. Laurer \parencite*[][p. 96-97]{laurer2024} points out the limitations of these models for text analysis, as they still rely on manually coded data to train and validate these models, but he also expresses the effectiveness for tasks such as classification, which previously would have been carried out manually across large corpuses. Similarly, the issue of multilingual corpuses has raised questions for practicality in research using machine automated analysis. Licht and Lind \parencite*{haukelicht2023} discuss multiple ways of addressing this problem, and view the use of multilingual corpuses as a way of broadening research scopes. Multilingual text analysis previously required multilingual analysts or multiple analysts, but recent innovation in machine translation has provided greater possibilities to approach multilingual topics \Parencite[][p. 4-6]{haukelicht2023}. One way of approaching multi-lingual corpuses is through input allignment, which involves translation of all texts into one language, often english, which can then be analysed just as a monolingual corpus \Parencite[][p. 10-11]{haukelicht2023}. Although this technique is not as accurate as human translation, it allows for large corpuses to be analysed, and remains largely accurate through various text-as-data techniques, but there do not seem to be any measures for testing the accuracy of manual versus automated multilingual analysis due to its complexity and specialisation \Parencite[][p. 20-21]{haukelicht2023}. 

Another important evolution in machine automated text analysis is the \textit{transformers revolution}. This comes from the training of various BERT --- a Google encoder based transformers model --- based models, one of which is BERTopic. BERTopic relies on clustering algorithms to determine texts that have similar word pairs and meaning in order to inductively create topics \Parencite[][p. 2-3]{grootendorst2022bertopic}. This has provided an increased ability to classify texts on a large scale for use in further topic modelling, and these clusterings can be saved for use on similar corpuses. Similar models have been developed for sentiment analysis as well, of which the Cardiff NLP Time LM has seen significant use \parencite*{camacho-collados-etal-2022-tweetnlp}. This roBERTa based model was trained on Twitter data, and has the highest accuracy across open-source models in recognizing overall sentiment, as well as particular sentiment patterns \Parencite[][p. 3-4]{camacho-collados-etal-2022-tweetnlp}. There are new transformers based models that are trained everyday for various tasks, but it is apparent that the use of these models will grow with the increasing amounts of data needed for text analysis in social sciences.

